\documentclass[11pt,a4paper]{article}
\usepackage[margin=24mm]{geometry}
\usepackage{fontspec}
\setmainfont{DejaVu Serif}[ItalicFont={DejaVu Serif},ItalicFeatures={FakeSlant=0.18},BoldItalicFont={DejaVu Serif Bold},BoldItalicFeatures={FakeSlant=0.18}]
\setmonofont{DejaVu Sans Mono}[Scale=0.8]
\usepackage{babel}
\babelprovide[main,import]{ukrainian}
\usepackage{amsmath,amssymb,booktabs,longtable,tabularx,graphicx,tikz,microtype}
\usetikzlibrary{arrows.meta,positioning}
\usepackage[unicode,hidelinks]{hyperref}
\setlength{\emergencystretch}{3em}
\title{Як побудувати мережу людей усередині міст\\для двомасштабного методу}
\author{Аналіз конструкції, даних та реалізації}
\date{16 жовтня 2026 року}
\begin{document}
\maketitle
\noindent У цьому варіанті вузол --- людина, спільнота --- місто її постійного проживання. Потрібно побудувати персональні контакти, а з них отримати локальні графи та міжміську матрицю методу. Нижче наведено конкретний проєкт побудови, перевірені джерела та обмеження. Це пояснювальний документ без анотації й висновків. Повну синтетичну країну та епідемічні траєкторії в межах цього аналізу не генеруємо.

\section{Основне рішення: одна людина, одна домашня спільнота}
Для кожної синтетичної людини $p$ фіксуємо місто проживання $h(p)=i$. Тоді
\[
V_i=\{p:h(p)=i\},\qquad N_i=|V_i|,\qquad V=\bigsqcup_iV_i.
\]
Потрібні також вік, домогосподарство, статус зайнятості та школа чи робоче місце. Людина може працювати в іншому місті, але її належність до $V_i$ протягом цього експерименту залишається домашньою. Зараження працівника під час роботи змінює його стан, а не домашнє місто.

Таким чином, локальне ребро з'єднує двох мешканців одного міста, навіть якщо вони зустрілися в іншому місці. Міжспільнотне ребро з'єднує мешканців різних міст, навіть якщо обидва працюють на одному підприємстві в третьому місті. Класифікувати контакт лише за місцем зустрічі було б помилкою для обраного поділу.

Це статична контактна модель зі сталим домашнім членством. Модель явного руху людей із щоденним розкладом і поверненням додому є іншим рівнем опису. Роботи про повторювану мобільність, зокрема Balcan--Vespignani, підкреслюють роль домашнього місця й часу повернення: \url{https://pmc.ncbi.nlm.nih.gov/articles/PMC3242939/}. Тут відвідування використовуватимемо для створення контактів, а не як готові події зараження.

\section{Які дані потрібні й що вже є}
У попередній підготовці збережено повну напрямлену таблицю INSEE 2022 про місця проживання та роботи, зведену до комун у географії 2025 року. Вона доступна як \path{../national_network_preparation/prepared/F_domestic.npz}; порядок комун --- у відповідному \path{nodes.csv}. Це корисний вхід призначення робочих місць. Матриця $W$ попереднього прикладу, де вузлами були комуни, не є $W$ нової персональної мережі.

Для наступної побудови джерела мають різні ролі:
\begin{itemize}
\item \href{https://www.insee.fr/fr/statistiques/8581696}{INSEE, Évolution et structure de la population en 2022}: чисельність і вікова структура, опубліковані в географії 01.01.2025. Потрібно брати рядки цілих комун, без додаткового повторного додавання округів Парижа, Ліона й Марселя.
\item \href{https://www.insee.fr/fr/statistiques/8647008}{INSEE, Couples--Familles--Ménages en 2022}: типи й склад домогосподарств. Це окремий потрібний набір; його детальну сумісність та всі комунальні обмеження ще треба перевірити перед генерацією домогосподарств.
\item \href{https://www.insee.fr/fr/statistiques/8582949}{INSEE, Mobilités professionnelles en 2022}: оцінені потоки зайнятих мешканців до місць роботи. Вони не дають адресатів персональних контактів, розмірів конкретних робочих груп чи мобільності дітей.
\item \href{https://www.insee.fr/fr/information/4802589}{INSEE, unités urbaines 2020, географія 2025}: альтернативний поділ на міські агломерації. Потрібно використовувати історичний випуск 2025, а не автоматично найновіший.
\item \href{https://journals.plos.org/plosone/article?id=10.1371/journal.pone.0133203}{Французьке дослідження контактів COMES-F}: підстава для вибору контактних характеристик за віком і середовищем. Конкретні контактні коефіцієнти в цьому аналізі не оцінено й не перенесено автоматично.
\end{itemize}

Методологічний прецедент --- синтетична французька популяція з домогосподарствами, школами й робочими місцями в дослідженні Faucher et al. (2022), \url{https://doi.org/10.1038/s41467-022-29015-y}. Загальні підходи до демографічно обґрунтованих контактних мереж описані в Mistry et al. (2021), \url{https://doi.org/10.1038/s41467-020-20544-y}. З цих робіт беремо принцип побудови, а не параметри COVID-19 чи твердження про точність нашої майбутньої мережі.

Виявлена пастка: компактний \href{https://www.insee.fr/fr/statistiques/8290591}{файл офіційної чисельності 2022} має межі комун 2024 року, тоді як наша мобільність --- 2025. Однаковий рік населення не гарантує однакової географії. Для кількісної перевірки далі використовуємо вікову базу 2022 у географії 2025; URL та хеші отриманих файлів будуть збережені окремо.

\section{Попередні побудови та готові набори}
Тут є прямі прецеденти, включно з масштабом країни. Водночас синтетична популяція з місцями діяльності та готовий персональний контактний граф --- різні продукти.

\paragraph{Франція, JRC, 2022.} У звіті \emph{Multipurpose synthetic population for policy applications}, розділ~4 та додаток~I, описано синтез 63~млн осіб для всієї Франції, домогосподарств і зв'язків із місцями роботи та навчання на основі INSEE. Це близький прецедент для нашого задуму. Проте зв'язок «людина--робоче місце» ще не визначає інтенсивності контакту між двома людьми. У пункті~A.5 звіт повідомляє про заплановане надання даних службам Єврокомісії через BDAP; це не підтвердження відкритого завантаження всього графа. Поточний публічний доступ до його бінарних даних у цьому аналізі не встановлено. Джерело: \url{https://publications.jrc.ec.europa.eu/repository/handle/JRC128595}.

\paragraph{США, Jiang et al., 2024.} Робота \emph{A Large-Scale Geographically Explicit Synthetic Population with Social Networks for the United States} надає, за розділом Data Records, 330\,526\,186 синтетичних осіб і домашні, робочі та навчальні мережі. Автори посилаються на дані за штатами в OSF; робочу мережу спочатку будували на національному рівні через міжштатні зв'язки. Це найпряміший кандидат для адаптації готового графа. Перш ніж його імпортувати, потрібно перевірити ID між файлами, збереження ребер через межі вибраної території та умови використання. Файли OSF тут ще не завантажено. Стаття: \url{https://doi.org/10.1038/s41597-024-03970-1}; дані: \url{https://doi.org/10.17605/OSF.IO/FPNC2}; код: \url{https://github.com/njiang8/geo-synthetic-pop-usa}.

\paragraph{Франція, Metz, Faucher et al., 2022.} Епідемічна агентна модель використовує синтетичну популяцію 117\,492 осіб із домогосподарствами, школами, робочими місцями та іншими середовищами контактів. Це приклад реалізації на рівні міста, корисний для перевірки деталей конструкції. Джерело: \url{https://doi.org/10.1038/s41467-022-29015-y}.

Отже, для дослідження саме похибки SMPT варто спочатку оцінити готовий американський набір на кількох сусідніх містах: це може прибрати значну частину роботи з генерації. Для французького сценарію вже маємо узгоджені комунальні потоки, а звіт JRC дає опублікований протокол синтезу. Потрібно окремо перевірити доступність його результату. Названі роботи підтверджують здійсненність конструкції, але не доводять точності нашого двомасштабного наближення.

\section{Що саме називати містом}
Адміністративна комуна зручна завдяки стабільному коду та прямому приєднанню INSEE. Проте вона може бути невеликим селом; велике фактичне місто може складатися з багатьох комун. Тому називати всі 34\,858 комун «містами» неправильно.

Для міського трактування придатна \emph{unité urbaine}: офіційна міська агломерація, що об'єднує відповідні комуни. Для нашої мети це вибір просторового поділу, а не спільнота, знайдена алгоритмом кластеризації ребер. Сільське населення все одно повинно належати рівно одній спільноті. Є два прозорі варіанти: залишити кожну сільську комуну окремо або ввести допоміжні сільські групи за департаментами. Другий варіант --- лише обчислювальний стресовий поділ: він може зібрати віддалені та незв'язні локальні компоненти. Усереднення $I_i/N_i$ тоді змішує їхні стани й може посилити похибку обходу мостових вершин. Для змістовного територіального сценарію слід залишити окремі комуни або обґрунтувати суміжні сільські кластери. Допоміжні групи департаментів не називатимемо містами.

В офіційній таблиці сільські комуни також мають значення \texttt{UU2020}, яке позначає збірну категорію поза міськими агломераціями. Групувати лише за цим кодом без поля типу означало б видати сільські території за місто. У збереженій таблиці використано явні префікси \texttt{U:} та \texttt{R:}.

Коди й межі фіксуємо в географії 01.01.2025. Перехід до агломерацій здійснюємо через таблицю членства комун. Для одного й того самого персонального графа можна порівняти кілька таких поділів, але не обирати поділ за найкращою вже побаченою точністю методу. Зокрема поділ, побудований із поїздок на роботу, сам по собі може поглинути велику частину міжкомунних зв'язків.

\section{Що показав фактичний аудит завантажених даних}
До наявних 34\,858 комун приєднано обидва нові джерела без втрати кодів і без дублювання рядків. У файлі населення 34\,903 рядки: крім цілих комун він містить муніципальні округи, які повторно до населення не додаємо. Таблиця міських одиниць має також комуни Майотти, що лежать поза обраним охопленням потоків. Усі наступні числа стосуються цього охоплення, тобто Франції без Майотти.

\paragraph{Відповідність кодів не означає повноту чисел.} Для Celles~(15031), Chalinargues~(15035), Chavagnac~(15047) і Sainte-Anastasie~(15171) рядки є, але населення й вікові поля мають значення NA. У поясненні INSEE до місцевих даних згадано відновлення цих комун 01.01.2025 та потребу розглядати їх разом із Neussargues-Moissac~(15141), залежно від вихідних меж статистики: \url{https://www.insee.fr/en/statistiques/6457611?geo=COM-43182}. У нашій таблиці пропуски залишено як NA. Сума доступних значень --- 67\,760\,573; вона не встановлює населення кожної з п'яти сучасних комун. Перед персональною генерацією потрібне узгодження цього спільного облікового контуру або документований розподіл із додаткових джерел. Присвоїти чотирьом NA нульове населення було б неправильно.

Для решти записів сума семи вікових груп збігається з загальною чисельністю до числової точності: найбільше абсолютне відхилення $6.5\cdot10^{-10}$. Водночас для 21 комуни оцінена кількість зайнятих мешканців у джерелі мобільності перевищує чисельність населення віком від 15 років у новій таблиці; найбільша різниця --- близько 346.45 особи. Це означає, що сирі маргінали цих наборів несумісні з прямим призначенням кожному працівникові особи віком від 15 років; це не автоматичне визнання одного джерела помилковим. Вона показує, де потрібне узгодження до призначення конкретних людей.

\paragraph{Розмір макрорівня залежить від обліку сільських територій.} В охопленні є 2457 міських агломерацій і 27\,307 сільських комун. Сукупність агломерацій і кожної сільської комуни окремо залишає майже 30 тисяч спільнот. Агломерації разом із допоміжними сільськими групами за департаментами дають 2553 спільноти. Населення найбільшої агломерації за доступними рядками --- 10\,944\,094; навіть один її локальний граф буде великим.

\begin{table}[htbp]
\centering\small
\input{analysis/v2/partition_table.tex}
\caption{Три поділи того самого охоплення. Внутрішні поїздки --- частка оцінених працівників у домашній та робочій спільноті, що збігаються; включено роботу у своїй комуні. Остання колонка --- лише нижня оцінка пам'яті п'яти щільних масивів $M\times M$ типу float64. Це не вимірювання часу, пікової пам'яті чи точності SMPT.}
\end{table}

Для кожного поділу обчислено $\Phi=P^\mathsf{T}FP$ через розріджену матрицю членства $P$ та перевірено збереження суми $F$. $\Phi$ описує поїздки між новими територіями й не підміняє контактну $W$. Зростання частки внутрішніх поїздок із 34.60\% до 65.38\% після укрупнення не є оцінкою похибки методу чи часткою локальної контактної ваги: контакти співпрацівників визначаються парами домашніх міст через спільне місце роботи.

\paragraph{Конкретний регіональний варіант.} Bretagne містить 1202 комуни й 3\,422\,845 мешканців за цим джерелом. У поточному $F$ близько 1\,372\,552 оцінених працівників одночасно живуть і працюють у регіоні, 45\,471 його мешканець працює в іншому внутрішньому регіоні, а 32\,500 мешканців інших регіонів працюють у Bretagne. Ще близько 1673 зайнятих мешканців мають місце роботи поза внутрішнім охопленням. Навіть цей менший варіант потребує явної політики меж і є завеликим для безумовного запуску поточного Full-коду.

Для масштабу доступної суми населення і гіпотетичного середнього ступеня~20 один CSR-масив сусідств із int32 ID та float64 вагами потребував би приблизно 15.65~GiB разом зі зсувами int64. Це арифметична оцінка $N\bar d(4+8)+8(N+1)$; вона не включає стани людей, шари, логи, копії чи Python-об'єкти й не є виміряним споживанням пам'яті нашою програмою.

\section{Від статистики до синтетичних людей}
Таблиця населення дає оцінені чисельності за віком, а не списки реальних осіб. Створюємо анонімні синтетичні записи, що відтворюють ці обмеження. Мінімальна схема: \texttt{person\_id}, \texttt{home\_commune}, \texttt{home\_city}, \texttt{age\_group}, \texttt{household\_id}, \texttt{activity\_type}, \texttt{activity\_place}, \texttt{contact\_group}.

Потрібно узгодити цілі числа $N_{ia}$ за комунами $i$ й віковими групами $a$. Окремо опубліковані округлені вікові поля можуть не сумуватися точно до загальної чисельності. Зберігаємо початкові значення, вибираємо пріоритетні маргінальні суми й робимо контрольоване цілочисельне узгодження. Звичайне незалежне округлення кожної клітинки не гарантує виконання всіх обмежень.

Потім формуємо домогосподарства за їхніми розмірами й складом, використовуючи відповідні таблиці або мікродані. Не варто незалежно жеребкувати віки й випадково ділити людей на сім'ї: це може створити демографічно неправдоподібні поєднання. IPF/IPU або зважування шаблонів домогосподарств допомагають узгодити розподіли, але потребують сумісних обмежень та перевірки здійсненності. Самі вікові суми не визначають єдиний склад домогосподарств.

Зайнятість призначаємо лише сумісним особам. Нехай $E_i$ --- кількість синтетичних зайнятих мешканців комуни $i$, а $n_{ir}$ --- ціла кількість із них, що працюють у місці $r$. Для повного охоплення призначень потрібно
\[
 n_{ir}\geq0,\quad n_{ir}\in\mathbb Z,\qquad \sum_r n_{ir}=E_i.
\]
У суму включаємо роботу вдома та поза вибраною територією, якщо вони є в джерелі. Початкова $F_{ir}$ містить оцінені, часто дробові числа працівників. Узгоджуємо її з доступними людьми й потрібними маргіналами, після чого робимо контрольоване округлення. Зберігаємо таблицю відхилень $n-F$, похибки сум і втрат малих потоків. Несумісний рядок не виправляємо мовчазним обрізанням: потрібно або переглянути зайнятість, або задокументувати корекцію даних.

Дітям та студентам потрібні окремі правила й джерела навчальної мобільності. Застосувати до них розподіл поїздок зайнятих дорослих означало б додати непідтверджену структуру. Поки немає цих даних, можна задати явний спрощений шкільний сценарій і не називати його відновленою реальною мобільністю.

\section{Чому поїздка між містами ще не є контактним ребром}
Позначимо $F_{ir}$ кількість мешканців $i$, що працюють у $r$. Це зв'язок \emph{походження--місце роботи}. Контактне ребро пов'язує двох людей. Працівники з міст $A$ та $B$ можуть зустрічатися на роботі в $C$, хоча ніхто з них не їздить із $A$ до $B$. Отже, ані $F$, ані $F+F^\mathsf{T}$ не є потрібною міжспільнотною матрицею персональної мережі.

Конкретний алгоритм: у кожному місці роботи формуємо підприємства, а в них контактні групи або команди. Призначаємо працівників із різних домашніх міст до цих груп. Розміри, галузь, віковий склад і спосіб змішування впливають на результат; одна OD-таблиця їх не ідентифікує. Один великий повний граф усіх працівників міста давав би необґрунтовану кількість контактів.

Для найпростішого контрольованого прикладу група $g$ розміру $m_g\geq2$ має $\kappa_g$ ефективних контактів на людину за день. Якщо всі пари в ній рівноправні, кожна отримує вагу
\[
 w_{pq}^{(g)}=\frac{\kappa_g}{m_g-1}.
\]
Тоді зважений ступінь кожного учасника цього шару дорівнює $\kappa_g$. Самотня людина не отримує ребра й не реалізує контактний бюджет групи. Для великих колективів використовуємо розріджені команди або інший явно описаний контактний механізм; наведена формула для кліки не є обов'язковою моделлю всіх підприємств.

\paragraph{Малий приклад.} На одному підприємстві в $C$ працюють четверо людей із домашніми містами $A,A,B,C$. При $\kappa=3$ усі шість ребер мають вагу~1. Отримуємо одне локальне ребро всередині $A$ і
\[
W_{AB}=2,\quad W_{AC}=2,\quad W_{BC}=1.
\]
Тут $W_{AB}>0$ навіть за $F_{AB}=F_{BA}=0$. Саме ці пари домашніх міст потрібні макрорівню. Приклад відтворено в \path{contact_algebra.py} та перевірено статичним тестом.

\section{Очікувані ваги можна оцінити до генерації графа}
Нехай $m_{irg}$ --- число працівників із дому $i$ у групі $g$ місця $r$. Для реалізованих груп-клік
\begin{equation}
W_{ij}^{\mathrm{work}}=
 \sum_{r,g:m_{rg}\geq2}\frac{\kappa_{rg}}{m_{rg}-1}\,m_{irg}m_{jrg},\qquad i\ne j.
\label{eq:realized}
\end{equation}
Це точна сума ваг між домашніми спільнотами для заданих груп. Коли домашні міста в межах кожного місця роботи перемішуються без переваг, розміри груп фіксовані, $T_r=\sum_i n_{ir}$, то
\[
\mathbb E[m_{irg}m_{jrg}]
=\frac{m_{rg}(m_{rg}-1)n_{ir}n_{jr}}{T_r(T_r-1)}.
\]
Звідси для $T_r\geq2$
\begin{equation}
\mathbb E[W_{ij}^{\mathrm{work}}]=
\sum_r\frac{n_{ir}n_{jr}}{T_r(T_r-1)}
\sum_{g:m_{rg}\geq2}\kappa_{rg}m_{rg}.
\label{eq:expected}
\end{equation}
Якщо всі групи мають принаймні двох людей і спільне $\kappa_r$, формула скорочується до $\sum_r\kappa_r n_{ir}n_{jr}/(T_r-1)$. Місця з менш ніж двома працівниками дають нульовий внесок. За сортування за віком, галуззю чи домашнім містом це спрощення не обґрунтоване.

Формула корисна для проєктування й перевірки генератора, але в порівнянні Full/SMPT потрібно підставляти $W$ фактично створеного графа. Підміна реалізованої матриці її математичним сподіванням додасть до похибки методу ще й невідповідність мереж.

\section{Зібрати шари й отримати вхід саме нашого методу}
Домогосподарства, класи, робочі команди й позадомашні контакти формуємо окремо. Для одного й того самого неорієнтованого ребра складаємо інтенсивності:
\[
w_{pq}=\sum_s\theta_s w_{pq}^{(s)}.
\]
Тут $\theta_s$ --- явно заданий коефіцієнт середовища. Поїздки вже використано для призначення людей; додавати ще одне ребро «за факт поїздки» без окремої контактної моделі означало б подвійний рахунок. Транспортний шар можна додати, якщо відомо або явно змодельовано спільний транспорт і час контакту.

Одиниці повинні бути однаковими: $w$ --- ефективні контакти за день, $\beta$ --- передача на одиницю контакту, $\beta w$ --- інтенсивність зараження за день. Це модельні інтенсивності, а не автоматичне перетворення кількості поїздок на ймовірність інфекції. Контактні параметри й відмінності середовищ ще потрібно обґрунтувати; у цьому документі їх не калібровано.

Після об'єднання ребер будуємо
\[
G_i=G[V_i],\qquad
W_{ij}=\sum_{p\in V_i}\sum_{q\in V_j}w_{pq},\quad i\ne j,
\qquad W_{ii}=0.
\]
Одне неорієнтоване зовнішнє ребро дає однакову вагу в $W_{ij}$ і $W_{ji}$. Це два можливі напрями передачі, а не подвоєння інтенсивності одного напряму. Внутрішні ребра залишаються в $G_i$. Незалежна контрольна рівність:
\[
\sum_i\sum_{e\in E(G_i)}w_e
+\frac12\sum_{i\ne j}W_{ij}=\sum_{e\in E(G)}w_e.
\]
Full Gillespie отримує $G$; двомасштабний метод --- ті самі $G_i$, $N_i$ й обчислену з них $W$. Зберігаємо також зовнішні сусідства людей, інакше неможливо коректно визначити граничні вершини й перевірити розміщення імпорту.

Базова макроінтенсивність $\beta W_{ij}(I_i/N_i)(S_j/N_j)$ усереднює стани кінців ребер. Реальна побудова сама по собі не робить це усереднення точним. Національні дані можуть створити дуже нерівномірні граничні ролі: більшість міжміських контактів нестимуть зайняті дорослі, а не випадковий мешканець міста.

У поточному API потрібно явно передати неоднакові розміри спільнот і нейтральні \texttt{alphas=[1,...,1]}, а параметри передачі Full і SMPT узгодити. Не можна покладатися на генератор однакових штучних спільнот або автоматичні спектральні корекції. Для базового порівняння окремо фіксуємо режим корекцій і фактор \texttt{T}, не допускаючи прихованого повторного масштабування ваг.

\section{Статичне усереднення, масштаб і межі території}
Мережа, що об'єднує дім і роботу за весь день, дозволяє обом шарам передавати інфекцію в будь-який момент модельного дня. Вона втрачає розклад і порядок контактів. Для оцінки похибки SMPT Full і SMPT спочатку мають використовувати однакову статичну мережу. Порівняння зі змінною в часі мережею --- окреме дослідження похибки часової агрегації.

При зменшенні населення в $c$ разів зберігаємо реалістичні розміри сімей і робочих команд та контактний бюджет людини, зменшуючи кількість груп. Тоді загальна вага контактів приблизно масштабується як $c$, а не $c^2$. Малий потік може округлитися до нуля й зникнути разом із маршрутом поширення. Тому масштабована країна --- окремий сценарій, для якого потрібно перевірити OD-похибки, зв'язність і граничні ролі. Ваги вершин $1/c$ не роблять такий CTMC автоматично еквівалентним повному населенню: змінюються флуктуації, розподіли прибуття та, для моделей з одужанням, імовірність згасання.

Для вибраного регіону або набору міст окремо зберігаємо потоки назовні й усередину з інших домашніх територій. Якщо залишити тільки місцевих жителів, зникнуть і їхні контакти з вилученими мешканцями, зокрема на роботі всередині регіону. Рішення --- розширити територію буфером або визначити зовнішню модель; мовчазне перенормування внутрішніх контактів не відновлює вилучену структуру. Із наявної таблиці не можна вивести всі імпорти з-за кордону: вона не містить повної сукупності осіб із зовнішнім місцем проживання.

\section{Що змінити в реалізації для масштабу країни}
Перевірка поточного коду показує, що сама наявність даних ще не означає можливості завантажити всю країну. У \path{two_layer_ctmc/engines/macro_engine.py} створюються щільні масиви розміру $M\times M$; розрахунок і вибір макроподії також обробляють щільні структури. Наведена в таблиці оцінка лише для п'яти масивів float64 є нижньою межею; код має додаткові масиви й тимчасові об'єкти.

На мікрорівні NetworkX та словники на кожне ребро мають суттєві накладні витрати. Оркестратор додатково зберігає компактні сусідства, а проби всередині кроку синхронізації копіюють локальні стани. Вибір сприйнятливого адресата та деякі операції з інфікованими переглядають списки людей. Логи всіх подій і щільні часові записи також можуть домінувати за пам'яттю. Тут немає вимірювання швидкодії чи заяви про прискорення.

Для повної країни потрібні розріджені макроребра з оновленням лише відповідних інтенсивностей, вибір події через дерево сум, масиви станів і CSR-сусідства людей, явне керування пробами й потоковими логами. Це окрема зміна симулятора. Як перший перевірюваний випадок доцільні кілька сусідніх міст із буфером або зменшена регіональна популяція, для якої доступний Full-еталон. Вибір першого випадку робимо за доступністю даних і ресурсів, а не за очікуваною малою похибкою.

\section{Варіант для контрольованої перевірки методу}
Якщо головне питання --- поведінка SMPT на багатьох містах, можна спершу побудувати географічно мотивований контрольований граф. Задаємо $N_i$, локальні контактні групи й цільову симетричну $W$. Для пари міст обираємо $K_{ij}\leq N_iN_j$ різних пар людей і присвоюємо кожній вагу $W_{ij}/K_{ij}$. Тоді сума реалізованих зовнішніх ваг точно дорівнює цільовій.

Можна змінювати концентрацію граничних контактів, частку зайнятих мостових осіб, повторне використання одних і тих самих кінців та зв'язок між внутрішнім і зовнішнім ступенями. Отримуємо контрольоване дослідження механізмів похибки. Це не відновлення реальної французької соціальної мережі, а окремий тип вхідних даних із точними конструктивними обмеженнями.

\section{Протокол наступної побудови та перевірки}
\begin{enumerate}
\item Зафіксувати країну, роки, географію, домашній поділ і політику зовнішньої території. Вирішити, чи адаптуємо готові американські дані, чи синтезуємо Францію з INSEE.
\item Перевірити ID, географічні відповідності, вікові й домогосподарські суми; підготувати цілочисельні демографічні та OD-обмеження з журналом корекцій.
\item Створити людей і домогосподарства; призначити заняття, місця роботи й навчання; сформувати контактні групи. Зберігати окремі випадкові потоки для цих етапів і незалежні від них seed-и епідемій.
\item Згенерувати контакти шарами, об'єднати повторні пари, відкинути самопетлі, перевірити скінченність і невід'ємність ваг. Зафіксувати одиниці й контрольні контактні характеристики.
\item Отримати $G_i$, реалізовану $W$, граничні вершини та їхні зовнішні ступені. Перевірити суму ваг незалежним блоковим підсумовуванням, унікальність людей і відсутність втрат сільського населення.
\item Зберегти \texttt{persons}, \texttt{groups}, \texttt{edges}, \texttt{communities}, \texttt{W}, джерела, хеші, параметри генератора, seed-и та відхилення від цільових статистик. Зміна коду або входів потребує нової папки результатів.
\item Лише після цього підготувати окремий затверджуваний епідемічний протокол. У цьому аналізі такі запуски не виконувалися.
\end{enumerate}

До оцінки точності методу потрібні щонайменше розподіли локальних зважених ступенів, частки ваг кожного шару, міжміські ваги, концентрація їхніх кінців, компоненти зв'язності й похибки демографічних та OD-маргіналів. Повна відповідність цих агрегатів не доводить правильності всіх персональних контактів: однакові $F$ та вікові суми допускають багато різних графів.

\section{Як ця конструкція допоможе розділити похибки SMPT}
Для фіксованого стану точна інтенсивність передачі між домашніми містами:
\[
a^*_{ij}=\beta\sum_{p\in V_i,q\in V_j}
w_{pq}\mathbf1\{X_p=I,X_q=S\}.
\]
Вона дозволяє окремо дослідити усереднення інтенсивності. Знання лише середніх станів граничних вершин не прибирає залежностей уздовж конкретних ребер. Для розміщення імпорту в місті $j$ точний умовний розподіл має вигляд
\[
\pi^*_{ij}(q)=
\frac{\mathbf1\{X_q=S\}\sum_{p\in V_i:X_p=I}w_{pq}}
{\sum_{u\in V_j:X_u=S}\sum_{p\in V_i:X_p=I}w_{pu}}.
\]
Ця формула визначена лише за наявності активного зовнішнього ребра $I$--$S$. Якщо наближений макрорівень пропонує подію, коли знаменник нульовий, потрібна явна політика несумісної події. Її відхилення вже використовує точну інформацію про інтенсивність і не є чистою корекцією лише адресата.

Для майбутнього порівняння розрізняємо базовий SMPT, варіант із граничними характеристиками, точні інтенсивності та точне розміщення. Поєднання обох точних компонентів відтворює міжспільнотний генератор Full лише за однакової локальної динаміки та точного узгодження подій у часі; скінченний крок синхронізації залишає окрему похибку.

Часи прибуття порівнюємо за однакового горизонту, початкового зараження та правила цензурування. Структурно недосяжні міста та неприбуття до кінця горизонту не зливаємо в одну вибірку скінченних часів. Реалістичність вихідного графа, похибка статичного усереднення та похибка SMPT потребують окремих оцінок. Жодної нової епідемічної точності чи часової переваги наведений аналіз не встановлює.

\section{Що збережено разом із цим документом}
\path{sources.json} містить URL, атрибуцію, статус джерел і SHA256. \path{fetch_inputs.py} перевіряє або отримує два офіційні архіви. \path{analyze_design.py} виконує детермінований аудит приєднання населення й міських одиниць до вже підготовлених потоків. У \path{analysis/v2/} збережено вихідні таблиці аудиту, підсумки й маніфест. \path{contact_algebra.py} та \path{test_design.py} задають і перевіряють малий приклад контактної конструкції. Це готові входи й специфікація наступного генератора; повного персонального графа Франції тут поки немає.

% NEXT_SECTIONS
\end{document}
